% !TEX root = chapter-standalone.tex

\section{Modules, Vector spaces}
\label{sec:vector-spaces}

\begin{ctdefinition}[Real vector space]
    \label{def:real-vector-space}
    A \maindef{real vector space} is
    \begin{body}
        \constit
        \begin{enumerate}
            \item a set~$\vecspBset$, the elements of which are called \emph{vectors};
            \item a binary operation~$+ \colon \vecspBset \cartprod \vecspBset \sto \vecspBset$, called \emph{vector addition};
            \item an element $0 \setin \vecspBset$;
            \item an operation ~$\cdot \colon \reals \cartprod \vecspBset \sto \vecspBset$, called \emph{scalar multiplication}.
        \end{enumerate}
        \condit
        \begin{enumerate}
            \item $\tup{\vecspBset, +}$ is a commutative group, with neutral element $0$;
            \item Scalar multiplication is a left action operation of the ring $\mathbb{R}$:
                  \begin{enumerate}
                      \item $(\lambda \mu) \cdot \ela = \lambda \cdot( \mu \cdot \ela) \quad \forall \lambda, \mu \setin \reals, \ \forall \ela \setin \vecspA$;
                      \item $1 \cdot \ela = \ela \quad \forall \ela \setin \vecspAset$;
                      \item $(\lambda + \mu) \cdot \ela = (\lambda \cdot \ela) + (\mu \cdot \ela) \quad \forall \lambda, \mu \setin \reals, \ \forall \ela \setin \vecspBset$;
                      \item $\lambda \cdot (\ela + \elb) = (\lambda \cdot \ela) + (\lambda \cdot \elb), \quad \forall \lambda \setin \reals, \ \forall \ela, \elb \setin \vecspBset$.
                  \end{enumerate}
        \end{enumerate}
    \end{body}
\end{ctdefinition}


    \subsection{Modules}

    \todotext{J: @J: explain modules briefly.... precisely the same notion as ring actions?}

\begin{ctdefinition}[Left ring action]
    \label{def:left-ring-action-operation}
    Let $\stylecat{R} = \tup{\stylesets{R}, *, +, 1, 0}$  be a ring and $\stylecat{A} = \tup{\stylesets{A}, \mathbf{+}_{\stylesets{A}}, \text{0}_\stylesets{A} }$ a commutative group. 
    A \maindef{left ring action-operation} of the ring $\stylecat{R}$ on the commutative group \stylecat{A} is 
    \begin{body}
        \constit
        \begin{enumerate}
            \item a function
            \begin{equation}
                  \label{eq:lact-ring}
                  \act \colon \makecartprod{\stylesets{R}, \stylesets{A}} \sto \stylesets{A},
              \end{equation}
              for which we use the infix-notation $\act(r, a) =  r \cdot a$ for all $r \in \stylesets{R}, a \in \stylesets{A}$. 
            \end{enumerate}
        \condit
        \begin{enumerate}
            \item $\act$ is a left action-operation of the monoid $\tup{\stylesets{R}, *, 1}$ on $\stylesets{A}$:
            \begin{enumerate}
                      \item $(\lambda * \mu) \cdot a = \lambda \cdot( \mu \cdot a) \quad \forall \lambda, \mu \setin \stylesets{R}, \ \forall a \setin \stylesets{A}$;
                      \item $1 \cdot a = a \quad \forall a \setin \stylesets{A}$;
            \end{enumerate}
            \item $\act$ is ``distributive'' with respect to the commutative groups $\tup{\stylesets{R}, +, 0}$ and $\tup{\stylesets{A}, \mathbf{+}_\stylesets{A}, \text{0}_\stylesets{A}}$,  respectively:
            \begin{enumerate}
                                        \item $(\lambda + \mu) \cdot \ela = (\lambda \cdot \ela) \mathbf{+}_\stylesets{A} (\mu \cdot \ela) \quad \forall \lambda, \mu \setin \stylesets{R}, \ \forall \ela \setin \stylesets{A}$;
                      \item $\lambda \cdot (\ela \mathbf{+}_\stylesets{A} \elb) = (\lambda \cdot \ela) \mathbf{+}_\stylesets{A} (\lambda \cdot \elb), \quad \forall \lambda \setin \stylesets{R}, \ \forall \ela, \elb \setin \stylesets{A}$.
                  \end{enumerate}
        \end{enumerate}
    \end{body}
\end{ctdefinition}



\todo{J: the typesetting of the sum with subscript above is off and needs, I guess, to be handled with a macro.}

\todotext{J: I'm starting to think that the speak of ``left action operation'' is too much of a mouthful, and it seems it would be enough to speak of a ``left action'' instead, as this can still then be differentiated from an ``action'', which is the simpler, cleaner definition simply in terms of a homomorphism.}

\begin{ctdefinition}
Let $\stylecat{R}$ be a ring. A left $\stylecat{R}$-module is a commutative group $\stylecat{A}$ together with a left ring action of $\stylecat{R}$ on $\stylecat{A}$.
\end{ctdefinition}


\begin{ctdefinition}[Ring action]
Let $\stylecat{R}$ be a ring and let $\stylecat{A}$ be a commutative group. A \emph{ring action} of $\stylecat{R}$ on $\stylecat{A}$ is a ring-homomorphism
\begin{equation}
                  \act \colon \stylecat{R} \mto \End(\stylecat{A}),
\end{equation}
where $\End(\stylecat{A})$ denotes the endomorphism ring of $\stylecat{A}$. 
\end{ctdefinition}

\subsection{Vector spaces}

\todotext{J: @J: add material}

\todotext{@J: define vector spaces}



\begin{example}
    $\tup{\reals^n, +, \cdot}$
\end{example}

\todotext{@J: expand $R^n$ example...}

\todotext{@J: example: polynomials (is there also a simple/neat applied interpretation/example here?)}

\begin{gradedexercise}[\exname{RealPolynomials}]
    \label{ex:real-polynomials}
    Let $\setA$ denote the set of polynomials in one variable and with coefficients in $\reals$.
    In other words, elements of $\setA$ are polynomials of the form
    \begin{equation}
        p(X) = a_nX^n + a_{n-1}
        X^{n-1} + \dots + a_1X + a_0,
    \end{equation}
    where $n \setin \natnumbers$ and $a_n, a_{n-1}, \dots, a_1, a_0 \setin \reals$ may vary.

    \begin{enumerate}
        \item Check (prove) that $\setA$ forms a commutative monoid when equipped with the usual addition operation for polynomials and the zero polynomial $p(X) = 0$ as neutral element.
        \item Is the commutative monoid $\tup{\setA, +, 0}$ in fact a group?
        \item Check that $\tup{\setA, +, 0}$ forms a real vector space when we equip it with the following (usual) scalar multiplication:
              \begin{equation}
                  \lambda \cdot (a_nX^n + a_{n-1}
                  X^{n-1} + \dots + a_1X + a_0) \definedas \lambda a_nX^n + \lambda a_{n-1}X^{n-1} + \dots + \lambda a_1X + \lambda a_0
              \end{equation}
              for any $\lambda \setin \reals$.
    \end{enumerate}
\end{gradedexercise}

\solutionof{RealPolynomials}

\todotext{@J: example: stuff related to fourier analysis}
